% Originally: new section (no predecessor in the week-based skeleton).
% Quelle: Linus Lüchinger/Slides/Slides-03-30.pdf, slides 12--13 (examples and
%         non-examples), 16, 23, 26 (one surface in three representations);
%         Linus Lüchinger/Slides/Slides-04-30.pdf, slide 7 (tangent-field test).

\section{Recognising a submanifold in practice}
\label{sec:recognising_submanifolds}

% Supplement: Linus Lüchinger/Slides/Slides-03-30.pdf, slides 16, 23, 26
% Why this section exists at all: rem:three_submanifold_descriptions says which description to
% reach for, but Corsin's notes never run one surface through all three, and give no worked
% non-example beyond the crossing in ex:cross_not_submanifold. Linus spends most of one session
% on exactly those two gaps, and this section follows his treatment. (Was an ainote until
% 2026-08-09; it is source-mining logistics, which Test 2 in style.md sends to a % comment.)

\subsection{One surface, three descriptions}

% Supplement: Linus Lüchinger/Slides/Slides-03-30.pdf, slides 16, 23, 26
The three characterisations of a submanifold are easiest to keep apart when they are applied to
the same object at the same point. Take the open upper half-sphere
\[\mathbb{S}^2_+ := \{(x,y,z) \in \mathbb{R}^3 : x^2+y^2+z^2 = 1,\ z > 0\},\]
a $2$-dimensional submanifold of $\mathbb{R}^3$, and throughout fix the point
\[p := \Big(0,\ \tfrac{1}{\sqrt2},\ \tfrac{1}{\sqrt2}\Big) \in \mathbb{S}^2_+.\]

\begin{example}[The upper half-sphere, three ways]
\label{ex:half_sphere_three_ways}
\begin{enumerate}[label=\textbf{(\alph*)}]
  \item \textbf{Implicit.} Put $F(x,y,z) := x^2+y^2+z^2-1$ on the open set
    $U := \{z>0\}$. Then $\mathbb{S}^2_+ = F^{-1}\{0\}$ and
    \[\nabla F(x,y,z) = (2x,\ 2y,\ 2z), \qquad \nabla F(p) = (0,\ \sqrt2,\ \sqrt2) \neq 0,\]
    so $\Jac F$ has full rank $\ell = 1$ everywhere on $U$ (indeed $2z > 0$ already forces the
    row to be nonzero). \cref{thm:regular_value_theorem} gives a submanifold of dimension
    $3-1 = 2$.
  \item \textbf{Graph.} Over the open unit disc $V := \{(x,y) : x^2+y^2 < 1\}$ put
    $f(x,y) := \sqrt{1-x^2-y^2}$. Then $\mathbb{S}^2_+ = \{(x,y,f(x,y)) : (x,y) \in V\}$ is the
    graph of $f$, and $f(0,1/\sqrt2) = \sqrt{1-1/2} = 1/\sqrt2$ recovers the third coordinate of
    $p$. Here the half-sphere is globally a graph, which is exactly why the \emph{upper half} is
    a friendlier example than the whole sphere.
  \item \textbf{Parametrization.} With $W := (0,\pi/2) \times (0,2\pi)$ and
    \[\phi(\theta,\psi) := (\sin\theta\cos\psi,\ \sin\theta\sin\psi,\ \cos\theta),\]
    we get $p = \phi(\pi/4,\ \pi/2)$. The Jacobian
    \[\Jac\phi(\theta,\psi) = \begin{pmatrix}
        \cos\theta\cos\psi & -\sin\theta\sin\psi \\
        \cos\theta\sin\psi & \sin\theta\cos\psi \\
        -\sin\theta & 0
      \end{pmatrix}\]
    is tall ($3\times2$), and full rank $=2$ means $D\phi$ is \emph{injective}
    (\cref{rem:full_rank_two_meanings}): the bottom row forces the first column to be nonzero
    whenever $\sin\theta \neq 0$, and the second column is then independent of it because
    $\sin\theta \neq 0$ makes $(-\sin\theta\sin\psi,\ \sin\theta\cos\psi)$ nonzero. On $W$ we
    have $\theta \in (0,\pi/2)$, so $\sin\theta > 0$ throughout.
\end{enumerate}
\end{example}

% Generator: Claude Opus 5 (High) --- FIG-17-04
% Geometry check. Screen projection of a 3D point (X,Y,Z) is (Y - 0.5X, Z - 0.3X), then
% scaled by the sphere radius R = 1.25. The marked point is p = (0, 1/sqrt2, 1/sqrt2),
% which has X = 0, so it projects to R*(0.7071, 0.7071) = (0.884, 0.884) -- a point at
% screen distance sqrt(0.884^2 + 0.884^2) = 1.250 = R from the origin, i.e. exactly on the
% dome's silhouette circle, at 45 degrees. That is why the dot sits ON the arc in all three
% panels, and why the gradient arrow leaves it along the outward radial direction (45 deg):
% nabla F(p) = (0, sqrt2, sqrt2) is radial, and radial directions in the y-z plane are
% undistorted by this projection.
\begin{center}
\begin{tikzpicture}[>=Latex, scale=0.95,
    dome/.style={blue!65!black, thick},
    equator/.style={blue!45!black, thin},
    panellabel/.style={font=\small, align=center}]

  \def\R{1.25}
  % All three captions sit on the SAME baseline y = -2.75, below every drawn element.
  \def\capy{-2.75}
  % --- reusable dome: outline arc + equator ellipse (back dashed, front solid) ---
  \newcommand{\halfsphere}{%
    \draw[dome] (-\R,0) arc[start angle=180, end angle=0, radius=\R];
    \draw[equator, dashed] (-\R,0) arc[start angle=180, end angle=0,
        x radius=\R, y radius={0.40*\R}];
    \draw[equator] (\R,0) arc[start angle=0, end angle=-180,
        x radius=\R, y radius={0.40*\R}];
  }

  % ============ Panel (a): implicit ============
  \begin{scope}[shift={(0,0)}]
    \halfsphere
    \coordinate (pa) at (0.884,0.884);
    \fill[black] (pa) circle (1.6pt);
    \node[font=\scriptsize, below left=-1pt and 0pt] at (pa) {$p$};
    \draw[orange!90!black, very thick, ->] (pa) -- ++(0.60,0.60);
    \node[orange!90!black, font=\scriptsize, above] at (1.30,1.46) {$\nabla F(p)$};
    \node[panellabel] at (0,\capy) {\textbf{(a)} implicit\\[1pt]
      {\footnotesize$\mathbb{S}^2_+ = F^{-1}\{0\}$}};
  \end{scope}

  % ============ Panel (b): graph ============
  \begin{scope}[shift={(4.6,0)}]
    \halfsphere
    \coordinate (pb) at (0.884,0.884);
    % The disc V, drawn detached below the dome so the drop line stays readable.
    \draw[green!45!black, thick, fill=green!12] (0,-1.30) ellipse ({\R} and {0.40*\R});
    \node[green!45!black, font=\scriptsize] at (-1.62,-1.30) {$V$};
    \fill[black] (pb) circle (1.6pt);
    \node[font=\scriptsize, below left=-1pt and 0pt] at (pb) {$p$};
    % Drop line from p to its shadow. p has (x,y) = (0, 1/sqrt2), whose screen abscissa is
    % R*0.7071 = 0.884 -- strictly inside the disc, since |(0,1/sqrt2)| = 0.707 < 1.
    \draw[gray, dashed] (0.884,0.884) -- (0.884,-1.30);
    \fill[green!45!black] (0.884,-1.30) circle (1.4pt);
    \node[font=\scriptsize, below right=-1pt and -2pt] at (0.884,-1.42) {$(x,y)$};
    \node[panellabel] at (0,\capy) {\textbf{(b)} graph\\[1pt]
      {\footnotesize$z = \sqrt{1-x^2-y^2}$}};
  \end{scope}

  % ============ Panel (c): parametrization ============
  \begin{scope}[shift={(9.2,0)}]
    \halfsphere
    \coordinate (pc) at (0.884,0.884);
    \fill[black] (pc) circle (1.6pt);
    \node[font=\scriptsize, below left=-1pt and 0pt] at (pc) {$p$};
    % Parameter rectangle W, drawn detached below the dome.
    \draw[orange!90!black, thick, fill=orange!12] (-1.15,-1.75) rectangle (0.15,-0.95);
    \node[orange!90!black, font=\scriptsize] at (-1.50,-1.35) {$W$};
    \fill[orange!90!black] (-0.50,-1.35) circle (1.4pt);
    \node[font=\scriptsize, below right=-1pt and -4pt] at (-0.50,-1.86)
      {$(\tfrac{\pi}{4},\tfrac{\pi}{2})$};
    \draw[->, thick, orange!90!black] (0.30,-1.15) to[out=15, in=255] (0.82,0.70);
    \node[orange!90!black, font=\scriptsize] at (1.15,-0.45) {$\phi$};
    \node[panellabel] at (0,\capy) {\textbf{(c)} parametrization\\[1pt]
      {\footnotesize$p = \phi(\tfrac{\pi}{4},\tfrac{\pi}{2})$}};
  \end{scope}
\end{tikzpicture}
\end{center}

% Generator: Claude Opus 5 (High)
\begin{remark}[Why all three theorems are stated \emph{locally}]
\label{rem:why_submanifolds_are_local}
The half-sphere is unusually obliging: one graph and one parametrization cover all of it. That is
not the general situation, and each description fails to globalise for its own reason.
\begin{itemize}
  \item \textbf{A parametrization need not be global.} The circle $\mathbb{S}^1$ is a
    $1$-dimensional submanifold, yet no single parametrization covers it. Suppose
    $\gamma : I \to \mathbb{S}^1$ were a homeomorphism from an open interval, as
    \cref{thm:local_parametrization} demands. Delete one point from each side: $I \setminus \{t\}$
    falls into \emph{two} pieces, while $\mathbb{S}^1 \setminus \{\gamma(t)\}$ is still an arc, so
    \emph{one} piece. A homeomorphism preserves the number of connected components
    (\cref{thm:continuity_preserves_connected}), so no such $\gamma$ exists. That is the counting
    argument of \cref{ex:cross_not_submanifold}, run in the opposite direction. Two overlapping
    arcs do cover the circle, and two is the minimum.
  \item \textbf{An implicit description need not be global either.} The Möbius strip is a
    $2$-dimensional submanifold of $\mathbb{R}^3$, but it is not $g^{-1}\{0\}$ for any single
    $g$ with $\nabla g \neq 0$: such a $g$ would supply a globally consistent normal direction
    $\nabla g/\lvert\nabla g\rvert$, and the Möbius strip has none
    (\cref{sec:orientable_manifolds}).
\end{itemize}
Requiring the descriptions only near each point costs nothing, since being a submanifold is
itself a local condition, and it buys all the examples that matter.
\end{remark}

\subsection{A gallery of non-examples}

% Supplement: Linus Lüchinger/Slides/Slides-03-30.pdf, slides 12--13
Linus prefaces the formal definitions with a list of sets that are \emph{not} submanifolds, which
is worth reproducing because each entry fails for a different reason. Being a $d$-dimensional
submanifold means looking, near every one of its points, like a flat piece of $\mathbb{R}^d$, and
there are several distinct ways to violate that.

\begin{itemize}
  \item \textbf{The figure eight} $\{(x,y) : y^2 = x^2 - x^4\}$ has a \emph{crossing}. Away from
    the origin it is a perfectly good curve. At the origin four branches meet. This is the same
    failure as \cref{ex:cross_not_submanifold}, and the same connectedness count settles it.
  \item \textbf{The lollipop} is a circle with a segment attached, or any \qt{line sticking out
    of a sphere}. Here the trouble is not a crossing but a \emph{change of dimension} at the
    junction: near the attachment point the set is neither $1$- nor $2$-dimensional.
  \item \textbf{The closed square} $[0,1]^2 \subseteq \mathbb{R}^2$ has a \emph{boundary}. Its
    interior is a fine $2$-dimensional submanifold, but no neighbourhood of a corner or edge
    point looks like an open piece of $\mathbb{R}^2$, because points of $[0,1]^2$ near an edge
    fill only a half-plane. (Sets like this are the subject of
    \cref{sec:submanifolds_with_boundary}; a submanifold \emph{with boundary} is a genuinely
    weaker notion, not a submanifold that happens to have edges.)
  \item \textbf{$\mathbb{Q}^n \subseteq \mathbb{R}^n$} has \emph{no dimension at all}. It is not
    $0$-dimensional, since a $0$-dimensional submanifold is a set of isolated points and no
    rational point is isolated. And it is not $n$-dimensional, since it contains no open ball.
\end{itemize}

% Generator: Claude Opus 5 (High) --- FIG-17-05
% Geometry check. Panel (a) is the lemniscate y^2 = x^2 - x^4, parametrized as
% (cos t, sin t cos t) for t in [0, 2pi); at t = pi/2 and t = 3pi/2 the curve passes through
% (0,0), which is why the self-crossing is at the origin and nowhere else. Panel (b): the
% circle has centre (-0.45,0) radius 0.62, so its rightmost point is (0.17, 0); the stick
% runs from (0.17,0) to (1.15,0), i.e. it starts exactly ON the circle -- the junction dot
% is at (0.17,0). Panel (c): the square is [0,1]^2 scaled to side 1.3; the marked edge point
% (0.65, 0) is the midpoint of the bottom edge, and the half-disc above it is the part of a
% small ball that actually lies in the square.
\begin{center}
\begin{tikzpicture}[>=Latex, scale=1.15,
    bad/.style={red!80!black},
    leader/.style={gray!70, very thin}]
  % All captions share the baseline y = -1.70; all annotations sit at y = -1.10,
  % clear of every drawn curve (the lemniscate reaches only |y| = 0.625).
  \def\capy{-1.70}
  \def\anny{-1.10}

  % ---- (a) figure eight: (x,y) = 1.25*(cos t, sin t cos t), which satisfies
  %      y^2 = x^2 - x^4 after undoing the scale. It passes through the origin
  %      exactly twice, at t = pi/2 and t = 3pi/2.
  \begin{scope}[shift={(0,0)}]
    \draw[blue!70!black, thick, domain=0:360, samples=200, smooth, variable=\t]
      plot ({1.25*cos(\t)}, {1.25*sin(\t)*cos(\t)});
    \fill[bad] (0,0) circle (1.8pt);
    \draw[leader] (0,\anny+0.16) -- (0,-0.14);
    \node[bad, font=\scriptsize] at (0,\anny) {crossing};
    \node[font=\small] at (0,\capy) {\textbf{(a)} figure eight};
  \end{scope}

  % ---- (b) lollipop: circle centred (-0.50,0) of radius 0.62, so its rightmost
  %      point is (0.12,0); the stick runs from there to (1.25,0), so the junction
  %      dot at (0.12,0) really does lie on both pieces.
  \begin{scope}[shift={(4.0,0)}]
    \draw[blue!70!black, thick] (-0.50,0) circle (0.62);
    \draw[blue!70!black, thick] (0.12,0) -- (1.25,0);
    \fill[bad] (0.12,0) circle (1.8pt);
    \draw[leader] (0.42,\anny+0.16) -- (0.16,-0.10);
    \node[bad, font=\scriptsize] at (0.42,\anny) {junction};
    \node[font=\small] at (0.30,\capy) {\textbf{(b)} lollipop};
  \end{scope}

  % ---- (c) closed square: [0,1]^2 drawn with side 1.3. The marked point is the
  %      midpoint of the bottom edge; the orange region is the part of a small ball
  %      about it that actually lies in the square -- the upper half only.
  \begin{scope}[shift={(7.9,0)}]
    \fill[blue!10] (0,-0.65) rectangle (1.3,0.65);
    \draw[blue!70!black, thick] (0,-0.65) rectangle (1.3,0.65);
    \fill[orange!30] (0.65,-0.65) -- ++(0.42,0)
      arc[start angle=0, end angle=180, radius=0.42] -- cycle;
    \draw[orange!90!black, thick] (0.65-0.42,-0.65)
      arc[start angle=180, end angle=0, radius=0.42];
    \fill[bad] (0.65,-0.65) circle (1.8pt);
    \draw[leader] (0.65,\anny+0.16) -- (0.65,-0.75);
    \node[bad, font=\scriptsize] at (0.65,\anny) {only half a ball};
    \node[font=\small] at (0.65,\capy) {\textbf{(c)} closed square};
  \end{scope}
\end{tikzpicture}
\end{center}

\subsection{Proving a negative: the tangent field test}

% Supplement: Linus Lüchinger/Slides/Slides-04-30.pdf, slide 7
\cref{ex:cross_not_submanifold} disproved a submanifold claim by counting connected components
after deleting a point. That argument is sharp but narrow: it needs the bad point to disconnect
the set. Linus gives a second, complementary criterion, which catches cases the counting argument
misses.

\begin{proposition}[Continuous tangent field along a curve]
\label{prop:continuous_tangent_field}
Let $M \subseteq \mathbb{R}^n$ be a $1$-dimensional $C^1$ submanifold. Then every point of $M$
has a neighbourhood on which there is a continuous map $T$ assigning to each $q$ a nonzero vector
$T(q) \in T_qM$.
\end{proposition}

\begin{proof}
Fix $p \in M$. By \cref{thm:local_parametrization} there are an open interval $I \subseteq
\mathbb{R}$, a neighbourhood $V$ of $p$ in $\mathbb{R}^n$, and a $C^1$ homeomorphism
$\gamma : I \to V \cap M$ with $D\gamma_t$ injective for every $t$. Injectivity of a linear map
$\mathbb{R} \to \mathbb{R}^n$ says exactly that $\gamma'(t) \neq 0$. Set
\[T(q) := \gamma'\big(\gamma^{-1}(q)\big), \qquad q \in V \cap M.\]
This is nonzero by the above, lies in $T_qM$ by \cref{prop:tangent_space_via_curves}, and is
continuous as a composition of the continuous maps $\gamma^{-1}$ and $\gamma'$.
\end{proof}

The test is used in contrapositive form: \textbf{if no continuous nonvanishing tangent field
exists near a point, the set is not a $1$-dimensional submanifold there.}

% Generator: Claude Opus 5 (High)
\begin{aiexample}[The lollipop, disproved by tangent fields]
\label{ex:lollipop_tangent_field}
Let $L$ be the lollipop of the gallery above and let $q$ be its junction point. Deleting $q$
leaves \emph{three} pieces where a curve would leave two, so the counting argument of
\cref{ex:cross_not_submanifold} already applies. But it is instructive to see the tangent-field
test reach the same verdict, since it localises the contradiction differently.

Suppose $L$ were a $1$-dimensional submanifold near $q$, and let $T$ be the tangent field
supplied by \cref{prop:continuous_tangent_field}. Approach $q$ along the stick: the tangent
direction there is constant, so $T$ tends to a vector along the stick. Approach $q$ along the
circle instead: the circle meets the stick at a right angle at $q$, so $T$ tends to a vector
perpendicular to the stick. A continuous map cannot have two different limits at the same point,
so no such $T$ exists.

\begin{remark}[When to reach for which]
Both tests disprove the same statement, and neither subsumes the other.
Counting components is the tool when the offending point \emph{separates} the set, and it needs
no differentiability. The tangent-field test is the tool when the set stays connected but the
direction of travel jumps. A cusp such as $\{y^2 = x^3\}$ at the origin is the standard case,
where removing the origin leaves the expected two pieces and only the tangent direction gives the
failure away.
\end{remark}
\end{aiexample}

\subsection{Transporting a submanifold by a diffeomorphism}

% Generator: Claude Opus 5 (High)
% Gap closed 2026-08-10: def:submanifold is stated through charts, and the fact that a
% diffeomorphism carries submanifolds to submanifolds of the same dimension follows from it in
% four lines, but it was nowhere in the document. ex:image_of_segment_submanifold below needs it,
% and so does any argument that moves a submanifold into better coordinates.
Each of the three descriptions above answers the question \qt{is this particular set a
submanifold?}. A different question arises as soon as a map is in play: given that $M$ is a
submanifold and $\Phi$ a diffeomorphism, what is $\Phi(M)$? Here the chart definition is the easy
one to use, and it is about the only place where that is so.

\begin{proposition}[Diffeomorphisms carry submanifolds to submanifolds]
\label{prop:diffeomorphism_preserves_submanifold}
Let $U, W \subseteq \mathbb{R}^n$ be open, let $\Phi : U \to W$ be a $C^k$-diffeomorphism, and let
$M \subseteq U$ be an $m$-dimensional $C^k$ submanifold of $\mathbb{R}^n$. Then $\Phi(M)$ is an
$m$-dimensional $C^k$ submanifold of $\mathbb{R}^n$. In particular the dimension is unchanged.
\end{proposition}

\begin{proof}
Fix $q \in \Phi(M)$ and write $q = \Phi(p)$ with $p \in M$. By \cref{def:submanifold} there are an
open neighbourhood $U_p \subseteq \mathbb{R}^n$ of $p$, an open $V \ni 0$, and a
$C^k$-diffeomorphism $\Psi : U_p \to V$ with $\Psi(p) = 0$ and
$\Psi(U_p \cap M) = (\mathbb{R}^m\times\{0\}) \cap V$. Replacing $U_p$ by $U_p \cap U$, we may
assume $U_p \subseteq U$.

Then $\Phi(U_p)$ is an open neighbourhood of $q$, and
\[\Psi \circ \Phi^{-1} : \Phi(U_p) \longrightarrow V\]
is a composition of $C^k$-diffeomorphisms, hence itself one, and it sends $q$ to $\Psi(p) = 0$.
Since $\Phi$ is injective, $\Phi(U_p) \cap \Phi(M) = \Phi(U_p \cap M)$, and therefore
\[(\Psi\circ\Phi^{-1})\big(\Phi(U_p) \cap \Phi(M)\big) = \Psi(U_p \cap M)
  = (\mathbb{R}^m\times\{0\}) \cap V .\]
That is a submanifold chart for $\Phi(M)$ at $q$, of the same dimension $m$. As $q$ was arbitrary,
$\Phi(M)$ is an $m$-dimensional $C^k$ submanifold.
\end{proof}

% Source: old_exams/fs2023/Probeprfg1.pdf (Probeprüfung Analysis II, undated), Aufgabe 7, p. 2
% Extractor: Claude Opus 5 (High)
% The paper names no lecturer; see old-exam-mining.md.
% Placed directly after prop:diffeomorphism_preserves_submanifold because it is the same question
% -- which operations preserve the property -- asked about four more of them. Parts (c) and (d)
% add products and transversal intersections, neither of which appears anywhere in this chapter,
% and (d) is the only place in the document where transversality is used.
\begin{exercise}[True or False \textnormal{(which operations preserve submanifolds?)}]
\label{ex:operations_preserving_submanifolds}
Decide whether each of the following is true or false, with justification.
\begin{enumerate}[label=\textbf{(\alph*)}]
  \item If $M$ is a $k$-dimensional submanifold of $\mathbb{R}^n$ and $N \subseteq M$ is open in
    $M$, then $N$ is itself a $k$-dimensional submanifold of $\mathbb{R}^n$.
  \item The union of two disjoint $k$-dimensional submanifolds of $\mathbb{R}^n$ is a
    $k$-dimensional submanifold of $\mathbb{R}^n$.
  \item The Cartesian product of a $k$-dimensional submanifold of $\mathbb{R}^m$ and an
    $l$-dimensional submanifold of $\mathbb{R}^n$ is a $(k+l)$-dimensional submanifold of
    $\mathbb{R}^{m+n}$.
  \item If $M$ and $N$ are $(n-1)$-dimensional submanifolds of $\mathbb{R}^n$ for $n \geq 2$, and
    $T_pM \neq T_pN$ for every $p \in M \cap N$, then $M \cap N$ is an $(n-2)$-dimensional
    submanifold of $\mathbb{R}^n$.
\end{enumerate}
\exinfo{This exercise is taken from the mock examination \emph{Probepr\"ufung Analysis II}
(undated) in \texttt{old\_exams/fs2023/}, Problem 7. That paper names no lecturer.}
\exsol{sol:operations_preserving_submanifolds}
\end{exercise}

\begin{remark}[Why the dimension cannot move]
\label{rem:dimension_is_a_diffeomorphism_invariant}
The proof transports one chart into another and never touches $m$, which is the whole content: a
diffeomorphism can bend a submanifold arbitrarily, but it cannot change how many independent
directions there are along it. This is worth stating because the tempting error runs the other
way, imagining that a map into $\mathbb{R}^2$ whose image \qt{spreads out} raises the dimension.
It cannot. A curve stays a curve.
\end{remark}

% Source: old_exams/examFS24.pdf (21 August 2024), Wahr/Falsch-Fragen, Aufgabe 6, p. 4
% Extractor: Claude Opus 5 (High)
% Split: parts 6.1-6.3 of this block are ex:local_diffeo_parameter_alpha in
% content/15-inverse-function-theorem/01-inverse-function-theorem.tex. The enumeration is continued
% at (d) with \setcounter so that the two halves read as one exam question, which is what they are.
\begin{exercise}[True or False \textnormal{(the image of a segment)}]
\label{ex:image_of_segment_submanifold}
Continuing \cref{ex:local_diffeo_parameter_alpha}: let $\Phi \in C^1(\mathbb{R}^2,\mathbb{R}^2)$
satisfy $\Phi(0,0) = (0,0)\transp$ and
$\Jac\Phi(0,0) = \left(\begin{smallmatrix} \alpha+5 & -3 \\ 3 & \alpha-5\end{smallmatrix}\right)$,
and take $\alpha = 0$. Decide whether each of the following is true or false, with justification.
\begin{enumerate}[label=\textbf{(\alph*)}]
  \setcounter{enumi}{3}
  \item For $r > 0$ small enough, the image of $\{(x,y) \in B_r(0,0) \mid y = 0\}$ under $\Phi$ is
    a $C^1$ submanifold of dimension $1$.
  \item For $r > 0$ small enough, the image of $\{(x,y) \in B_r(0,0) \mid y = 0\}$ under $\Phi$ is
    a $C^1$ submanifold of dimension $2$.
\end{enumerate}
\exinfo{This exercise is taken from the exam of 21 August 2024 (Prof.\ Joaquim Serra), where it appears as the fourth and
fifth parts of the sixth block of true/false questions. Its first three parts are
\cref{ex:local_diffeo_parameter_alpha}.}
\exsol{sol:image_of_segment_submanifold}
\end{exercise}

% Source: exercises_2024/mock.pdf (Analysis II mock examination, Spring Semester 2024,
% Prof. Joaquim Serra), Exercise 9, p. 2
% Extractor: Claude Opus 5 (High)
% The R^3 -> R^2 counterpart of ex:local_diffeo_parameter_alpha / ex:image_of_segment_submanifold
% above, and not a duplicate of them: there the map is square and the question is invertibility,
% here it is wide and the question is what a surjective differential buys. Parts (a) and (d) are
% the ones with teeth --- the curve gamma passes through the point but does not lie in M, which
% both of them turn on.
% Part (4) of the printed block carries no "True / False" prompt, unlike its five siblings;
% transcribed as a statement to be decided, in line with the rest.
\begin{exercise}[True or False \textnormal{(a level set in $\mathbb{R}^3$, and a curve through it)}]
\label{ex:level_set_of_a_wide_jacobian}
Let $F \in C^1(\mathbb{R}^3,\mathbb{R}^2)$ satisfy
\[F(0,0,1) = \begin{pmatrix} 2 \\ -2\end{pmatrix} \qquad\text{and}\qquad
  \Jac F(0,0,1) = \begin{pmatrix} 1 & 1 & -2 \\ 0 & 1 & -1\end{pmatrix},\]
and set $M := F^{-1}\big(\{(2,-2)\transp\}\big)$. For $t < 1$ let
\[\gamma(t) := \Big(\sin(2t),\ t+\tfrac12t^2,\ \tfrac{\cos t}{1-t}\Big).\]
Decide whether each of the following is true or false, with justification.
\begin{enumerate}[label=\textbf{(\alph*)}]
  \item $(F\circ\gamma)'(0) = (0,0)\transp$.
  \item In a small neighbourhood of $(0,0,1)$, the set $M$ is a smooth submanifold of
    dimension $2$.
  \item In a small neighbourhood of $(0,0,1)$, the set $M$ is a smooth submanifold of
    dimension $1$.
  \item The vector $\gamma'(0)$ is normal to $M$ at the point $(0,0,1)$.
  \item There exist $\delta > 0$ and $\varphi,\psi \in C^1((-\delta,\delta),\mathbb{R})$ such
    that $F(\varphi(y),\,y,\,\psi(y)) = (2,-2)\transp$ for all $y \in (-\delta,\delta)$.
  \item For all sufficiently small $\varepsilon > 0$, the set
    $F\big(B_\varepsilon((0,0,1))\big)$ is open in $\mathbb{R}^2$.
\end{enumerate}
\exinfo{This exercise is taken from the mock examination for Analysis II, Spring Semester 2024
(Prof.\ Joaquim Serra), where it appears as the ninth block of multiple-choice questions.}
\exsol{sol:level_set_of_a_wide_jacobian}
\end{exercise}

% Transition: Claude Opus 5 (High)
Every argument in this section leaned on knowing which vectors point \emph{along} $M$ at a given
point. The next section makes that space of directions an object in its own right.
